% 第 3 章 §3.3（印刷页 127–133）：扩张 TQFT 与高阶范畴
% 对应 工程文件/work/en/ch03.txt 第 1001–1263 行（§3.3 标题至 §3.3.5 末句）。
% 注意：第 951–1000 行是 §3.2 末尾（定理 3.2.8 及其讨论，同在印刷页 127），不属本段。
% 本段无编号公式、无插图（图 3.1、图 3.2 在 §3.4，见 ch03.txt 第 1534–1599 行）。

\section{扩张 TQFT 与高阶范畴}

Atiyah--Witten 定义的一个基本局限，是把边界限制在余维 $1$。这固然使故事干净简单，却使我们无法通过把流形切成简单小块来计算 TQFT；$D=1,2$ 的情形相当特殊，而下一个更高维的分类定理——$D=3$ 时属于 Reshetikhin--Turaev [25] 的定理——则需要用到圆圈与余维 $2$ 的切割。

直观上，扩张 TQFT（extended TQFT）是从范畴 $\mathsf{Bord}^{\mathrm{or}}_D$ 出发的函子，只是被增强了足够多的结构，使得可以把流形切成带任意余维角的简单小块。这样讲故事多少是倒着来的；而且很快就会发现，我们需要把配边范畴中的两层结构——$(D-1)$ 维的对象与 $D$ 维的态射——换成 $(D+1)$ 层结构，一直下降到点。

\subsection{高阶范畴}

把这样的代数结构——在对象之上叠 $D$ 层态射、允许多维复合——编码起来的各种方式，统称为 $D$-范畴（$D$-category）。一个典型的例子是配边 $D$-范畴 $\mathsf{Bord}^{\mathrm{or}}_D$，其中（定向的）$0$ 维流形是对象，带边界的定向 $1$ 维流形是 $0$ 维流形之间的态射，带角的 $2$ 维流形是「$2$-态射」，如此类推。人们很快发现，这套游戏的代数规则已不再那么明确（事实上，可以设想出许多套规则），所以下面给出第一个归纳式的定义。

\begin{definition}
严格 $D$-范畴（strict $D$-category）是一个范畴，其中所有态射集 $\Hom(x,y)$ 都具有 $(D-1)$-范畴的结构，而且复合 $\Hom(x,y)\times\Hom(y,z)\to\Hom(x,z)$ 都是 $(D-1)$-范畴的严格双函子。复合严格结合。\footnote{因此「相等」确实指相等，而不是典范同构。} 两个 $D$-范畴之间的函子 $\varphi$ 是指底层范畴之间的函子，使得在态射上诱导的映射 $\varphi_*\colon\Hom(x,y)\to\Hom(\varphi x,\varphi y)$ 都是 $(D-1)$-范畴的函子。
\end{definition}

我们关心的是（$\C$-）线性范畴，它们在最高的两层上重现了向量空间与线性映射。习惯上还要求这些范畴容许有限直极限，至少如此（任意余极限有时更方便）。在 $D$-范畴的世界里，对任何 $D$ 都存在一个「单位对象」；所谓单位，是说一切东西都是它上面的模，正如每个 $\C$-向量空间都是 $\C$ 上的模。在线性范畴中，这个单位就是有限维向量空间范畴 $\mathsf{Vect}$。（在 $\C$-线性范畴中，每个对象都可以与一个有限维向量空间作张量积，而这对对象与向量空间双方都是函子性的。）再往上是线性范畴、线性函子与自然变换构成的严格 $2$-范畴，然后是严格 $2$-范畴构成的严格 $3$-范畴，依此类推。若假定直极限存在，则要求线性函子右正合，即保持直极限。

\begin{remark}
另一条熟悉的、把「向量空间与线性映射的范畴」提升一步的途径（而不是走上线性范畴的 $2$-范畴），是由代数、双模与交织映射（intertwining map）构成的 $2$-范畴 $\mathsf{Alg}$。可以把后者「全忠实」地嵌入线性范畴：把每个代数送到它的模范畴。（关键观察是：对两个代数 $A,B$，任何从 $A$-模到 $B$-模的右正合函子，都必然由在 $A$ 上与某个 $B$-$A$-双模作张量积诱导。）这是一个非严格的 $2$-范畴：态射的复合——双模的张量积——只在相差自然同构的意义下才有定义。
\end{remark}

\subsection{严格范畴与弱范畴}

处理范畴的经验应当提示严格定义中的缺陷：一旦身处范畴之中，我们该要求「性态良好的同构」，而不是相等。何谓性态良好，取决于手头的问题，例如在某些运算下的函子性。举例来说，要求复合运算严格结合似乎就不对——我们应当要求一个结合子（associator）$\alpha$，即由不同次序复合所得到的、从 $\Hom(x,y)\times\Hom(y,z)\times\Hom(z,w)$ 到 $\Hom(x,w)$ 的两个函子之间的自然同构。对这个结合子有一个自然条件（Stasheff 五边形恒等式）；确认它之后，我们就可以按「结合性在相干的结合子意义下成立」来工作，正如我们对待严格结合性那样。类似地，我们可以只要求 $2$-范畴之间的函子把 $1$-态射的复合保持到相干地选定的 $2$-同构为止。随着范畴深度的推进，可以放松为「相干的」同构的恒等式越来越多。把结合与交换的数据、以及它们之上的相干条件系统地罗列出来，就得到弱 $D$-范畴（weak $D$-category）的理论。当前高阶范畴的许多工作，动机正是要找出既有章法、又方便实用的方式来编码这些弱数据及其相干条件；见 [4, 19, 26] 以及众多其他文献。

直觉并不完全可靠。对 $D=2$ 结果没有问题：任何弱 $2$-范畴都可以「严格化」；但再往上就会遇到麻烦，我们将在下文一个群胚的例子中看到。暂时搁下这个麻烦，先来讨论一个简单的二维有限群规范理论。

\subsection{二维的有限群规范理论}

在下一节讨论一般的对偶理论之前，这个例子会有用。我们构造一个从无定向配边 $2$-范畴 $\mathsf{Bord}_2$ 到 $\mathsf{Alg}$ 的 $2$-函子。

与点对应，我们取群代数 $A:=\C\langle F\rangle$，即线性组合 $\sum_{f\in F}a_f\cdot e_f$（$a_f\in\C$），其乘法 $e_f\cdot e_{f'}=e_{ff'}$ 模仿群乘法。$A$-模同于 $F$ 的复表示。注意 $A\cong A^{\mathrm{op}}$，同构由反对合（anti-involution）$f\leftrightarrow f^{-1}$ 给出；于是比方说，$A$-$A$-双模与 $A^{\otimes 2}$-模之间没有区别。（这不适用于规范理论的扭转版本，见下文注 3.3.6；这也与如下事实有关：有限群的无扭转规范理论不需要流形上的定向。）

现在与任何区间对应，我们取双模 $A$；按端点的情形——两端都出、两端都入、或一出一入——把它看作左 $A^{\otimes 2}$-模、右 $A^{\otimes 2}$-模或 $A$-$A$-双模。换言之，这个区间可以是从 $\varnothing$ 到 $\{*,*\}$ 的配边、从 $\{*,*\}$ 到 $\varnothing$ 的配边，或是 $\{*\}$ 上的恒等态射。（如前所示，我们不需要给点 $*$ 定定向。）注意自然同构 $A\otimes_A A\cong A$，它表达的正是复合 $\Id\circ\Id=\Id$。

与闭圆周对应，我们取空间 $A\otimes_{A\otimes A}A$。把圆周呈现为复合 $\subset\circ\supset$，就迫使这个计算成立；它把 $Z_F(S^1)$ 等同于 $F$ 上、在 $F$ 的共轭作用下的余不变函数空间。后者自然同构于 $F$ 上的类函数空间，正好重现了 §3.1.2 中讨论的规范理论的向量空间。

下面两层既已定义，现在来考虑带角的曲面。要把它读成一个 $2$-态射，我们必须补充更多数据：把各个角标记为源或靶，并相合地把各条边（$1$-态射）标记为由该曲面定义的 $2$-态射的源或靶。§3.1.3 中那样的扭转规范理论情形还要糟糕得多，是指数级的：我们得为所有这些对象配上相合的定向。把全部可能性与相应的映射一一列出，很快就变得繁冗不堪。更好的构造又一次来自物理学家的路径积分。

\begin{annote}
底本此处作「the case of twisted gauge theories as in Subsection 3.1.3」，但底本的 §3.1.3 是「初步分类，$D=1$」；扭转情形见 §3.1.2 中编号为 3.1.3 的注「扭转（Twistings）」（底本印刷页 111）。译文照排底本的「§3.1.3」。[来源：底本 p.111 的注 3.1.3 与小节 3.1.3 标题对照；本译本 ch03a.tex 中该注编号亦为 3.1.3]
\end{annote}

于是设 $\Sigma$ 是带角的曲面，读作其边界的一个入部分与一个出部分之间的 $2$-态射。这些边界部分 $\partial_\pm\Sigma$ 是源与靶都相同的 $1$-态射。也就是说，$\partial_-\Sigma$ 上的角必须由边界的某一部分 $B$ 与 $\partial_+\Sigma$ 上的角相配，$B$ 则充当对象上的恒等。\footnote{说来容易混淆：在开–闭 TQFT 的文献中，这些边界区间的叫法不一，或称自由（free），或称受约束（constrained）。} 现在我们规定：与一个 $1$-态射（带边界的 $1$ 维流形）对应的空间 $A$，应解释为在边界上平凡化的 $F$-丛的群胚 $\mathrm{Bun}_F$ 的 $H^0$。类似地，运算 $Z_F(\Sigma)$ 由熟知的对应图上的 $(\pi_+)_*\circ\pi_-^*$ 在 $H^0$ 上给出：

\begin{center}
\begin{tikzcd}[row sep=1.6em, column sep=2.6em]
& \mathrm{Bun}_F(\Sigma,B) \arrow[dl, "\pi_+"'] \arrow[dr, "\pi_-"] & \\
\mathrm{Bun}_F(\partial_+\Sigma,B_+) & & \mathrm{Bun}_F(\partial_-\Sigma,B_-)
\end{tikzcd}
\end{center}

\subsection{对应 \texorpdfstring{$2$-范畴}{2-范畴}}

我们刚才所做的，是把规范 TQFT 从 $\mathsf{Bord}_2\to\mathsf{Alg}$ 经由有限群胚的对应（correspondence）所构成的 $2$-范畴 $\mathsf{Cor}_2$ 作分解。$\mathsf{Cor}_2$ 中的对象是有限群胚 $X_\bullet$；这个记号应当理解为包含对象集 $X_0$、态射集 $X_1$、源与靶映射 $X_1\rightrightarrows X_0$、恒等截面 $X_0\to X_1$，以及复合律。群胚之间的映射 $X_\bullet\to Y_\bullet$ 是对象上与态射上的一对映射，与这些结构相合。

$\mathsf{Cor}_2$ 中的 $1$-态射是群胚之间的对应 $Y_\bullet\leftarrow X_\bullet\to Z_\bullet$，而 $X_\bullet$ 与 $X'_\bullet$ 之间的 $2$-态射仍是形如 $X_\bullet\leftarrow C_\bullet\to X'_\bullet$ 的对应，且要与到 $Y_\bullet\times Z_\bullet$ 的投影相合。$1$-态射与 $2$-态射的复合是同伦纤维积：即把其中一个映射换成纤维化之后所作的朴素纤维积——所谓纤维化，是指这样的态射：一旦找到源或靶的一个提升，每条箭头就都能提升。

\begin{example}
点就是 $*\rightrightarrows*$。群胚 $BF$（$F$ 为有限群）描述为 $F\rightrightarrows*$。映射 $*\to BF$ 不是纤维化，但把 $*\rightrightarrows*$ 换成等价的群胚 $F\times F\rightrightarrows F$（它表示 $F$ 在自身上的平移作用）之后，就可以转成纤维化。于是同伦纤维积 $*\times_{BF}*$ 是 $F^3\rightrightarrows F^2$，它等价于集合 $F$（没有非平凡的态射）。试验证之！
\end{example}

\begin{remark}
这个弱 $2$-范畴是可严格化的，因为群胚的每个态射都可以函子性地转成纤维化。
\end{remark}

把主 $F$-丛的群胚赋给流形，就得到一个从 $\mathsf{Bord}_2$ 到 $\mathsf{Cor}_2$ 的 $2$-函子。可以把它看作取值于 $\mathsf{Cor}_2$ 的量子场论。接下来的任务，是定义一个从 $\mathsf{Cor}_2$ 到 $\mathsf{Alg}$ 的 $2$-函子，我们如下进行：

\begin{enumerate}[label=(\roman*),leftmargin=3em,itemsep=1pt]
\item 对 $\mathsf{Cor}_2$ 的每个对象，赋以该群胚的路径代数（path algebra）；
\item 对一个对应 $Y_\bullet\leftarrow X_\bullet\to Z_\bullet$，赋以 $Y_0\times Z_0$ 的所有对象上的直和，其求和项是 $X$ 在该处的同伦纤维上的函数；这是 $Y_\bullet\times Z_\bullet$ 的路径代数上的一个模；
\item 对一个由对应构成的对应 $C_\bullet$，用推–拉构造在函数空间之间定义一个线性映射。（联系两个函数的矩阵系数，计数同时映入两者的点，再按相对于第二个映射的自同构加权（即除以自同构群的阶）；这在有限规范理论中已经熟悉。）
\end{enumerate}

对 (ii) 还有一个更具概念的描述：若 $f\colon X_\bullet\to Y_\bullet$ 是群胚的纤维化，则常值层的直接像 $f_*\C$ 是 $X_\bullet$ 上的平坦向量丛——即一个带有箭头的可复合提升的丛。它在 $X_0$ 上的截面于是给出 $Y_\bullet$ 的路径代数上的一个模。这里应当想象流形的淹没（submersion）$f\colon X\to Y$：沿纤维的上同调 $R^if_*\C$ 都是带有典范平坦联络的向量丛。

\begin{exercise}
证明 $BF$ 的路径代数是 $\C\langle F\rangle$。从映射 $*\to BF$ 出发，证明这样得到 $F$ 的正则表示。对 $G\subset F$ 及诱导映射 $BG\to BF$，证明所得到的是诱导表示，即 $F/G$ 上的函数。
\end{exercise}

\begin{remark}[扭转规范理论]
一个类 $[\tau]\in H^2(BF;\C^\times)$ 定义了 $F$ 被 $\C^\times$ 的中心扩张（central extension）。我们把中心扩张设想为群上的一条线丛，全空间上带有相合的乘法：即同构 $\alpha_{f,g}\colon L_f\otimes L_g\to L_{fg}$，满足下方方块交换所表达的相干性：

\begin{center}
\begin{tikzcd}[row sep=2.4em, column sep=3.6em]
L_f\otimes L_g\otimes L_h \arrow[r, "{\Id\otimes\alpha_{g,h}}"] \arrow[d, "{\alpha_{f,g}\otimes\Id}"'] & L_f\otimes L_{gh} \arrow[d, "{\alpha_{f,gh}}"] \\
L_{fg}\otimes L_h \arrow[r, "{\alpha_{fg,h}}"'] & L_{fgh}
\end{tikzcd}
\end{center}

当然，所有直线都可以与 $\C$ 等同，此时 $\alpha_{f,g}$ 成为一个 $\C^\times$-值的群上闭链。改变同构 $L_h\cong\C$ 会使这个上闭链改变一个上边缘。

扭转群代数 ${}^\tau\C\langle F\rangle$ 定义为这条线丛的截面空间，它带有显然的乘法，提升群元素的乘积。它的模是 $F$ 的 $\tau$-射影表示。不过这一次，反代数是 ${}^{-\tau}\C\langle F\rangle$。扭转规范理论定义在定向流形上，但不经过无定向配边范畴分解。

建议读者对定向流形构造这个扭转规范理论；在检验各映射的相合性时，曲面相对于其边界的基本链，以及取值于 $\C^\times$ 的上同调的 Stokes 定理，都应当派上用场。
\end{remark}

\subsection{严格范畴的不足}

我们提到过，高阶范畴的严格定义虽然看似容易，却不合适。把注意力限制到高阶群胚（higher groupoid）——即所有态射都可逆的范畴——就能指认出一个确切的问题。人们假定，任何合理的高阶群胚理论都等价于拓扑中的同伦型（homotopy type）理论；具体地说，$D$-群胚应当对应于 $D$-型（$D$-type），即维数 $D$ 以上同伦群为零的拓扑空间。（宣布映射的同伦是一个等价关系，相当于在高阶范畴中宣布所有态射可逆。）于是，对任何好的 $D$-范畴定义，限制到群胚都应当产生同伦 $D$-型。然而，严格归纳定义在 $D=2$ 以上就有毛病：

\begin{proposition}
一个连通的同伦型 $X$ 能由严格群胚表示，当且仅当它在 $k_2$ 之后的 $k$-不变量全为零。这又成立当且仅当 $X$ 是某个群的分类空间，而该群是 $\pi_1(X)$ 被一个拓扑阿贝尔群所作的扩张。

特别地，若 $X$ 单连通，则它必须等价于一个拓扑阿贝尔群的分类空间；也就是若干 Eilenberg--MacLane 空间的乘积。
\end{proposition}

大多数单连通的同伦型都不是这种形状：确实，限制到这类空间几乎抹掉了同伦论中所有有趣的部分。例如，从我们前面讨论过的 $2$-球面出发，杀掉 $3$ 维以上的同伦群、只保留 $\pi_2=\pi_3=\Z$，便得到一个同伦 $3$-型。（这可以靠不断粘接维数递增的胞腔做到。）但我们已经看到，所得结果与乘积 $3$-型 $\CP^\infty\times K(\Z;3)$ 不同：只需用从 $\CP^2$ 到这两个空间的映射来检验便知。

该定理的证明模仿拓扑空间 $\pi_2$ 交换性的证明：关键的引理是，一个带有两个严格交换的结合乘法的空间，实际上严格交换，而且这两个乘法必须一致。在严格 $3$-群胚的世界里，这证明了二次环路空间的严格交换性，并推出 Whitehead 括号的消失。我们简要回顾一下这个故事。
